\documentclass[11pt,a4paper]{article}

%------------------------------------------------
% PACKAGES
%------------------------------------------------

\usepackage[utf8]{inputenc}
\usepackage[margin=1in]{geometry}
\usepackage{amsmath,amssymb,amsfonts}
\usepackage{graphicx}
\usepackage{booktabs}
\usepackage{hyperref}
\usepackage{xcolor}
\usepackage{float}
\usepackage{cite}
\usepackage{enumitem}
\usepackage{titlesec}
\usepackage{fancyhdr}
\usepackage{listings}

%------------------------------------------------
% COLORS
%------------------------------------------------

\definecolor{navyblue}{RGB}{0,32,96}

%------------------------------------------------
% PAGE STYLE
%------------------------------------------------

\setlength{\headheight}{14pt}

\pagestyle{fancy}
\fancyhf{}

\rhead{VisionForge Engineering Report}
\lhead{CS480: Computer Vision Engineering}
\rfoot{Page \thepage}

%------------------------------------------------
% HYPERLINK SETTINGS
%------------------------------------------------

\hypersetup{
    colorlinks=true,
    linkcolor=navyblue,
    citecolor=navyblue,
    urlcolor=navyblue
}

%------------------------------------------------
% DOCUMENT
%------------------------------------------------

\begin{document}

%================================================
% CUSTOM TITLE PAGE
%================================================

\begin{titlepage}
    \centering

    \vspace*{1cm}

    {\Huge \bfseries VisionForge \par}

    \vspace{0.4cm}

    {\Large
    \textbf{
    A Modular Classical Computer Vision Architecture for
    Multistage Image Processing, Spatial Feature Extraction,
    and Video Scene Dynamics
    }
    \par}

    \vspace{0.6cm}

    \rule{\linewidth}{1.5pt}

    \vspace{1.5cm}

    {\Large \textbf{Coursework Project Engineering Report} \par}

    \vspace{0.3cm}

    {\large Course: CS480 Computer Vision Engineering \par}

    \vspace{3cm}

    \begin{minipage}{0.8\textwidth}
        \centering
        \large

        \textbf{Prepared By:}\\[0.2cm]

        \textbf{Parth Jangir}\\

        Registration No.: \textbf{24BAI10175}\\[0.3cm]

        Department of Computer Science \& Engineering\\

        School of Computer Science and Engineering
    \end{minipage}

    \vfill

    {\large \textbf{Submission Date:} September 2026 \par}

    \vspace{1cm}

\end{titlepage}

%------------------------------------------------
% START ARABIC PAGE NUMBERING
%------------------------------------------------

\newpage
\pagenumbering{arabic}

%================================================
% ABSTRACT
%================================================

\begin{abstract}

VisionForge is a modular classical computer vision architecture designed to demonstrate the complete processing lifecycle of a computer vision system, from low-level image preprocessing to spatial feature extraction, segmentation, geometric analysis, pattern classification, and video scene dynamics.

The system is organized as a five-stage pipeline. The first stage performs image preprocessing and frequency-domain analysis using grayscale conversion, Gaussian and median filtering, contrast enhancement, histogram analysis, and the two-dimensional Discrete Fourier Transform (DFT). The second stage extracts multi-scale visual features using Canny edge detection, Laplacian of Gaussian (LoG), Difference of Gaussians (DoG), Harris corner detection, Gabor filters, and Histogram of Oriented Gradients (HOG). The third stage performs image segmentation using K-Means clustering, Mean-Shift clustering, region growing, and marker-based watershed segmentation.

The fourth stage focuses on geometric and structural analysis through Hough line detection, Hough circle detection, contour extraction, shape descriptors, Principal Component Analysis (PCA), and K-Nearest Neighbours (KNN). The fifth stage extends the architecture to video analysis using background subtraction, dense optical flow, sparse Lucas--Kanade optical flow, and affine motion estimation.

For a $512 \times 512$ RGB input image, the complete static processing pipeline achieved an observed execution time of approximately 8.3015 seconds. The benchmark produced 364 Harris corners, 142,884 HOG descriptors, 23 Hough lines, 31 circles, and 28 closed contours. The measurements demonstrate how different classical computer vision algorithms contribute to overall computational latency.

The project emphasizes modularity, interpretability, mathematical foundations, and measurable system performance rather than reliance on deep learning. The resulting architecture provides a practical framework for understanding how classical computer vision techniques can be combined into a complete engineering pipeline.

\end{abstract}

%================================================
% 1. INTRODUCTION
%================================================

\section{Introduction and Theoretical Motivation}

Computer vision is the field of artificial intelligence concerned with enabling computers to acquire, process, analyse, and interpret visual information. Although modern computer vision systems frequently rely on deep neural networks, classical computer vision remains important because many fundamental image understanding operations can be implemented using mathematically interpretable and computationally efficient algorithms.

VisionForge was developed to demonstrate how a collection of classical computer vision techniques can be integrated into a modular processing architecture. Instead of treating image processing, feature extraction, segmentation, shape detection, and motion analysis as isolated topics, the project combines them into a structured pipeline.

The architecture consists of five major processing stages:

\begin{enumerate}
    \item Low-level image preprocessing and frequency-domain analysis.
    \item Multi-scale spatial feature extraction.
    \item Image segmentation and region analysis.
    \item Geometric shape detection and pattern classification.
    \item Video scene dynamics and motion estimation.
\end{enumerate}

The overall objective is to understand not only the mathematical principles behind individual algorithms but also their practical computational behaviour when combined into an engineering system.

%================================================
% 2. SYSTEM ARCHITECTURE
%================================================

\section{System Architecture}

The VisionForge architecture follows a modular pipeline in which each stage performs a specific visual analysis task.

\begin{figure}[H]
    \centering
    \fbox{
        \begin{minipage}{0.9\textwidth}
            \centering
            \textbf{Input Image / Video}\\
            $\downarrow$\\
            \textbf{Preprocessing and Frequency Analysis}\\
            $\downarrow$\\
            \textbf{Feature Extraction}\\
            $\downarrow$\\
            \textbf{Segmentation}\\
            $\downarrow$\\
            \textbf{Shape and Pattern Analysis}\\
            $\downarrow$\\
            \textbf{Motion and Scene Dynamics}\\
            $\downarrow$\\
            \textbf{Visual Information / Measurements}
        \end{minipage}
    }
    \caption{Overall VisionForge processing architecture}
    \label{fig:architecture}
\end{figure}

The modular design allows individual algorithms to be tested independently while also enabling their outputs to be used as inputs for subsequent processing stages.

%================================================
% 3. IMAGE PREPROCESSING
%================================================

\section{Low-Level Image Preprocessing and Frequency-Domain Analysis}

Low-level preprocessing is the first stage of the VisionForge pipeline. The purpose of this stage is to improve image quality, reduce noise, normalize contrast, and extract frequency-domain information.

%------------------------------------------------
% 3.1 Grayscale
%------------------------------------------------

\subsection{Grayscale Conversion}

RGB images contain three colour channels. For many computer vision algorithms, processing all three channels is unnecessary. Therefore, the input image is converted into grayscale.

The grayscale intensity is calculated using the standard luminance relationship:

\begin{equation}
Y = 0.299R + 0.587G + 0.114B
\end{equation}

where $R$, $G$, and $B$ represent the red, green, and blue channels respectively.

The grayscale representation reduces the dimensionality of the input while preserving important intensity information.

%------------------------------------------------
% 3.2 Gaussian Filtering
%------------------------------------------------

\subsection{Gaussian Filtering}

Gaussian filtering is used to smooth an image and reduce high-frequency noise.

The two-dimensional Gaussian function is:

\begin{equation}
G(x,y)=
\frac{1}{2\pi\sigma^2}
\exp
\left(
-\frac{x^2+y^2}{2\sigma^2}
\right)
\end{equation}

where $\sigma$ controls the amount of smoothing.

The filtered image is obtained through convolution:

\begin{equation}
I_{G}(x,y)=I(x,y)*G(x,y)
\end{equation}

Gaussian filtering is particularly useful before edge detection because it suppresses small intensity fluctuations that could otherwise produce false edges.

%------------------------------------------------
% 3.3 Median Filtering
%------------------------------------------------

\subsection{Median Filtering}

Median filtering is a nonlinear filtering technique commonly used to remove impulse noise.

For each pixel, a neighbourhood is considered and the pixel value is replaced by the median of the neighbourhood values.

Unlike linear averaging filters, median filtering preserves edges relatively well while removing isolated noisy pixels.

%------------------------------------------------
% 3.4 CLAHE
%------------------------------------------------

\subsection{Contrast Limited Adaptive Histogram Equalization}

Contrast Limited Adaptive Histogram Equalization (CLAHE) improves local image contrast.

Instead of applying histogram equalization to the entire image, CLAHE divides the image into small regions called tiles and performs local histogram equalization.

A contrast limit is applied to prevent excessive amplification of noise.

The process can be summarized as:

\begin{enumerate}
    \item Divide the image into local tiles.
    \item Calculate the histogram of each tile.
    \item Clip histogram bins exceeding the contrast limit.
    \item Redistribute the clipped pixels.
    \item Apply histogram equalization.
    \item Interpolate between neighbouring tiles.
\end{enumerate}

The measured execution time for CLAHE was approximately 0.0890 seconds.

%------------------------------------------------
% 3.5 Histogram Analysis
%------------------------------------------------

\subsection{Histogram Analysis}

An image histogram represents the frequency of different intensity values.

For a grayscale image, the histogram can be represented as:

\begin{equation}
H(k)=\sum_{x}\sum_{y}[I(x,y)=k]
\end{equation}

where $k$ represents an intensity level.

The histogram provides information about brightness and contrast distribution.

The measured histogram analysis stage required approximately 0.0903 seconds.

%------------------------------------------------
% 3.6 DFT
%------------------------------------------------

\subsection{Two-Dimensional Discrete Fourier Transform}

The two-dimensional Discrete Fourier Transform converts an image from the spatial domain into the frequency domain.

The DFT is defined as:

\begin{equation}
F(u,v)=
\sum_{x=0}^{M-1}
\sum_{y=0}^{N-1}
f(x,y)
e^{-j2\pi
\left(
\frac{ux}{M}+\frac{vy}{N}
\right)}
\end{equation}

where $f(x,y)$ represents the spatial-domain image and $F(u,v)$ represents its frequency-domain representation.

The magnitude spectrum can be calculated as:

\begin{equation}
|F(u,v)| =
\sqrt{
\operatorname{Re}(F)^2+
\operatorname{Im}(F)^2
}
\end{equation}

The dominant measured frequency in the benchmark was approximately 9 cycles/pixel.

The FFT processing stage required approximately 0.4107 seconds.

%================================================
% 4. FEATURE EXTRACTION
%================================================

\section{Multi-Scale Feature Extraction}

Feature extraction converts an image into meaningful representations that describe edges, corners, textures, shapes, and local gradients.

VisionForge implements several classical feature extraction techniques.

%------------------------------------------------
% 4.1 Canny
%------------------------------------------------

\subsection{Canny Edge Detection}

Canny edge detection identifies locations of significant intensity changes.

The algorithm consists of several steps:

\begin{enumerate}
    \item Gaussian smoothing.
    \item Gradient computation.
    \item Gradient magnitude and direction calculation.
    \item Non-maximum suppression.
    \item Double thresholding.
    \item Edge tracking by hysteresis.
\end{enumerate}

The gradient magnitude can be represented as:

\begin{equation}
G=\sqrt{G_x^2+G_y^2}
\end{equation}

The measured Canny edge density was approximately 0.0300.

%------------------------------------------------
% 4.2 LoG
%------------------------------------------------

\subsection{Laplacian of Gaussian}

The Laplacian of Gaussian combines Gaussian smoothing with the Laplacian operator.

It can be expressed as:

\begin{equation}
LoG=\nabla^2(G_{\sigma}*I)
\end{equation}

The LoG operator is useful for detecting rapid intensity changes and blob-like structures.

The measured execution time was approximately 0.0091 seconds.

%------------------------------------------------
% 4.3 DoG
%------------------------------------------------

\subsection{Difference of Gaussians}

Difference of Gaussians approximates scale-space filtering by subtracting two Gaussian-blurred versions of an image.

\begin{equation}
DoG(x,y)=
G(x,y,k\sigma)*I-
G(x,y,\sigma)*I
\end{equation}

The measured execution time was approximately 0.0039 seconds.

%------------------------------------------------
% 4.4 Harris
%------------------------------------------------

\subsection{Harris Corner Detection}

Harris corner detection identifies points where image intensity changes significantly in multiple directions.

The structure tensor is:

\begin{equation}
M=
\begin{bmatrix}
\sum I_x^2 & \sum I_xI_y\\
\sum I_xI_y & \sum I_y^2
\end{bmatrix}
\end{equation}

The Harris response is:

\begin{equation}
R=\det(M)-k(\operatorname{trace}(M))^2
\end{equation}

where $k$ is an empirical constant.

The benchmark detected 364 Harris corners.

The execution time was approximately 0.0041 seconds.

%------------------------------------------------
% 4.5 Gabor
%------------------------------------------------

\subsection{Gabor Filters}

Gabor filters are useful for texture analysis.

A two-dimensional Gabor function can be represented as:

\begin{equation}
g(x,y)=
\exp
\left(
-\frac{x'^2+\gamma^2y'^2}{2\sigma^2}
\right)
\cos
\left(
2\pi\frac{x'}{\lambda}+\psi
\right)
\end{equation}

where

\begin{equation}
x'=x\cos\theta+y\sin\theta
\end{equation}

and

\begin{equation}
y'=-x\sin\theta+y\cos\theta
\end{equation}

Gabor filtering was one of the more computationally expensive feature extraction operations, requiring approximately 0.6864 seconds.

%------------------------------------------------
% 4.6 HOG
%------------------------------------------------

\subsection{Histogram of Oriented Gradients}

Histogram of Oriented Gradients (HOG) represents an image using local gradient orientations.

The gradient magnitude is:

\begin{equation}
m(x,y)=
\sqrt{
G_x(x,y)^2+G_y(x,y)^2
}
\end{equation}

The gradient orientation is:

\begin{equation}
\theta(x,y)=
\tan^{-1}
\left(
\frac{G_y(x,y)}{G_x(x,y)}
\right)
\end{equation}

The benchmark generated 142,884 HOG descriptors.

The measured HOG execution time was approximately 0.2715 seconds.

%================================================
% 5. SEGMENTATION
%================================================

\section{Spatial and Color Segmentation}

Segmentation divides an image into meaningful regions based on similarity in colour, intensity, texture, or spatial properties.

%------------------------------------------------
% 5.1 K-Means
%------------------------------------------------

\subsection{K-Means Clustering}

K-Means partitions data into $K$ clusters.

The objective function is:

\begin{equation}
J=
\sum_{i=1}^{n}
\sum_{j=1}^{K}
\mathbf{1}(c_i=j)
\|x_i-\mu_j\|^2
\end{equation}

where $x_i$ is a data point and $\mu_j$ is the centroid of cluster $j$.

The main steps are:

\begin{enumerate}
    \item Select the number of clusters $K$.
    \item Initialize cluster centroids.
    \item Assign each point to the nearest centroid.
    \item Recalculate centroids.
    \item Repeat until convergence.
\end{enumerate}

For $K=5$, the measured execution time was approximately 2.0897 seconds.

This was the largest individual measured computational cost in the static pipeline.

%------------------------------------------------
% 5.2 Mean Shift
%------------------------------------------------

\subsection{Mean-Shift Segmentation}

Mean-Shift is a non-parametric clustering method that searches for modes in a feature distribution.

The algorithm repeatedly shifts a candidate point toward the mean of neighbouring data points until convergence.

Unlike K-Means, Mean-Shift does not require the number of clusters to be specified explicitly.

%------------------------------------------------
% 5.3 Region Growing
%------------------------------------------------

\subsection{Region Growing}

Region growing begins from one or more seed points and expands a region by adding neighbouring pixels that satisfy a similarity criterion.

The general procedure is:

\begin{enumerate}
    \item Select seed pixels.
    \item Examine neighbouring pixels.
    \item Compare neighbouring pixels with the region criterion.
    \item Add suitable pixels to the region.
    \item Continue until no suitable pixels remain.
\end{enumerate}

%------------------------------------------------
% 5.4 Watershed
%------------------------------------------------

\subsection{Marker-Based Watershed Segmentation}

Watershed segmentation treats an image as a topographic surface.

Bright regions can be interpreted as peaks while dark regions represent valleys.

Marker-based watershed uses predefined markers to control the segmentation process and reduce unwanted over-segmentation.

%================================================
% 6. SHAPE AND PATTERN ANALYSIS
%================================================

\section{Geometric Shape Detection and Pattern Classification}

This stage extracts geometric information from segmented or edge-based images.

%------------------------------------------------
% 6.1 Hough Lines
%------------------------------------------------

\subsection{Hough Line Transform}

The Hough Transform can detect lines in an image.

A line can be represented using normal form:

\begin{equation}
\rho=x\cos\theta+y\sin\theta
\end{equation}

where $\rho$ represents the distance from the origin and $\theta$ represents the angle of the normal.

The benchmark detected 23 lines.

The measured execution time was approximately 0.0172 seconds.

%------------------------------------------------
% 6.2 Hough Circles
%------------------------------------------------

\subsection{Hough Circle Detection}

A circle can be represented as:

\begin{equation}
(x-a)^2+(y-b)^2=r^2
\end{equation}

where $(a,b)$ is the centre and $r$ is the radius.

The Hough Circle Transform searches for combinations of centre coordinates and radius values that provide strong evidence of circular structures.

The benchmark detected 31 circles.

The measured execution time was approximately 0.0160 seconds.

%------------------------------------------------
% 6.3 Contours
%------------------------------------------------

\subsection{Contour Extraction}

Contours represent continuous boundaries of objects.

Contour extraction can be used to identify object boundaries and calculate geometric properties.

The benchmark detected 28 closed contours.

The measured contour-processing time was approximately 0.0040 seconds.

%------------------------------------------------
% 6.4 Shape Descriptors
%------------------------------------------------

\subsection{Shape Descriptors}

Two useful geometric descriptors are circularity and solidity.

Circularity is defined as:

\begin{equation}
C=
\frac{4\pi A}{P^2}
\end{equation}

where $A$ is area and $P$ is perimeter.

Solidity is defined as:

\begin{equation}
S=
\frac{A}{A_{\text{convex hull}}}
\end{equation}

These descriptors allow geometric characteristics of objects to be compared.

%------------------------------------------------
% 6.5 PCA
%------------------------------------------------

\subsection{Principal Component Analysis}

Principal Component Analysis (PCA) is used for dimensionality reduction and identifying directions of maximum variance.

The covariance matrix is calculated as:

\begin{equation}
\Sigma=
\frac{1}{n-1}
X^TX
\end{equation}

Eigenvectors of the covariance matrix define principal directions, while the corresponding eigenvalues represent the amount of variance along those directions.

%------------------------------------------------
% 6.6 KNN
%------------------------------------------------

\subsection{K-Nearest Neighbours}

KNN classifies an unknown sample based on the classes of its nearest training samples.

For Euclidean distance:

\begin{equation}
d(x,y)=
\sqrt{
\sum_{i=1}^{n}
(x_i-y_i)^2
}
\end{equation}

The class is generally selected according to the majority vote among the $K$ nearest neighbours.

%================================================
% 7. VIDEO SCENE DYNAMICS
%================================================

\section{Video Scene Dynamics}

The final module extends the classical computer vision pipeline from static image processing to video analysis.

The main objective is to identify moving objects and estimate motion between consecutive video frames.

%------------------------------------------------
% 7.1 Background Subtraction
%------------------------------------------------

\subsection{Background Subtraction}

Background subtraction separates foreground objects from a relatively stationary background.

VisionForge uses the MOG2 background subtraction approach.

The basic process is:

\begin{enumerate}
    \item Capture the current frame.
    \item Estimate the background model.
    \item Compare the current frame with the background.
    \item Identify pixels belonging to moving objects.
    \item Generate the foreground mask.
\end{enumerate}

%------------------------------------------------
% 7.2 Dense Optical Flow
%------------------------------------------------

\subsection{Farnebäck Dense Optical Flow}

Dense optical flow estimates motion vectors for a large number of pixels between consecutive frames.

The resulting motion field provides information about the direction and magnitude of scene movement.

Dense optical flow is useful when motion information is required across a large portion of the image.

%------------------------------------------------
% 7.3 Sparse Optical Flow
%------------------------------------------------

\subsection{Lucas--Kanade Sparse Optical Flow}

The Lucas--Kanade method estimates motion for selected feature points.

It is based on the brightness constancy assumption:

\begin{equation}
I(x,y,t)=I(x+\Delta x,y+\Delta y,t+\Delta t)
\end{equation}

Using a first-order Taylor approximation leads to the optical flow constraint equation:

\begin{equation}
I_xu+I_yv+I_t=0
\end{equation}

where $u$ and $v$ represent the horizontal and vertical velocity components.

%------------------------------------------------
% 7.4 Affine Motion
%------------------------------------------------

\subsection{Affine Motion Estimation}

Affine transformation models translation, rotation, scaling, and shear.

A two-dimensional affine transformation can be written as:

\begin{equation}
\begin{bmatrix}
x'\\
y'
\end{bmatrix}
=
\begin{bmatrix}
a & b\\
c & d
\end{bmatrix}
\begin{bmatrix}
x\\
y
\end{bmatrix}
+
\begin{bmatrix}
t_x\\
t_y
\end{bmatrix}
\end{equation}

Affine motion estimation is useful for describing global camera or object motion.

%================================================
% 8. EXPERIMENTAL SETUP
%================================================

\section{Experimental Setup}

The VisionForge system was evaluated using a $512\times512$ RGB image.

The experiment measured execution time for individual operations and recorded the outputs produced by different algorithms.

The benchmark was performed on the complete static image-processing pipeline.

The measured results are summarized below.

%================================================
% 9. PERFORMANCE BENCHMARK
%================================================

\section{Empirical Latency Evaluation}

\begin{table}[H]
\centering
\caption{Measured execution time of VisionForge processing stages}
\label{tab:benchmark}
\begin{tabular}{lr}
\toprule
\textbf{Processing Operation} & \textbf{Time (s)}\\
\midrule
Grayscale Conversion & 0.0020\\
Gaussian Filtering & 0.0010\\
Histogram Equalization & 0.0000\\
CLAHE & 0.0890\\
Histogram Analysis & 0.0903\\
FFT / DFT & 0.4107\\
Canny Edge Detection & 0.0020\\
Laplacian of Gaussian & 0.0091\\
Difference of Gaussians & 0.0039\\
Harris Corner Detection & 0.0041\\
HOG & 0.2715\\
Scale-Space Processing & 0.4078\\
Gabor Filtering & 0.6864\\
K-Means ($K=5$) & 2.0897\\
Hough Line Detection & 0.0172\\
Hough Circle Detection & 0.0160\\
Contour Processing & 0.0040\\
\midrule
\textbf{Total Pipeline} & \textbf{8.3015}\\
\bottomrule
\end{tabular}
\end{table}

The total measured execution time of the static pipeline was 8.3015 seconds.

%================================================
% 10. FEATURE SUMMARY
%================================================

\section{Feature and Detection Summary}

The major outputs obtained during the benchmark are summarized below.

\begin{table}[H]
\centering
\caption{VisionForge benchmark feature summary}
\label{tab:features}
\begin{tabular}{lr}
\toprule
\textbf{Feature / Measurement} & \textbf{Value}\\
\midrule
Image Mean Intensity & 82.37\\
Image Standard Deviation & 61.81\\
Dominant FFT Frequency & 9 cycles/pixel\\
Canny Edge Density & 0.0300\\
Harris Corners & 364\\
HOG Descriptor Dimension & 142,884\\
Detected Hough Lines & 23\\
Detected Hough Circles & 31\\
Closed Contours & 28\\
\bottomrule
\end{tabular}
\end{table}

These measurements demonstrate that the pipeline produces both numerical image statistics and higher-level structural information.

%================================================
% 11. BOTTLENECK ANALYSIS
%================================================

\section{Computational Bottleneck Analysis}

The latency measurements show that different algorithms contribute significantly different amounts of computational overhead.

The K-Means segmentation stage required approximately 2.0897 seconds, representing approximately 25.17\% of the total measured pipeline execution time.

The Gabor filtering stage required approximately 0.6864 seconds, making it another significant contributor to computational cost.

The FFT and scale-space processing stages also produced noticeable overhead, requiring approximately 0.4107 seconds and 0.4078 seconds respectively.

The remaining operations generally required considerably less time.

These results demonstrate why profiling is important when constructing a computer vision system. An algorithm that is mathematically simple may still become computationally expensive when applied repeatedly to large images or multiple feature scales.

%================================================
% 12. PRACTICAL OBSERVATIONS
%================================================

\section{Practical Observations}

Several observations were obtained during the development and benchmarking process.

\subsection{Modularity}

Separating the system into independent modules makes it easier to test, debug, replace, and optimize individual algorithms.

\subsection{Lighting Sensitivity}

Classical image processing techniques can be sensitive to lighting conditions. Contrast enhancement techniques such as CLAHE can improve local visibility, but excessive enhancement can also amplify noise.

\subsection{Computational Cost}

The benchmark demonstrates that segmentation and texture analysis can become computational bottlenecks.

\subsection{Interpretability}

Unlike many black-box models, classical computer vision algorithms provide mathematically interpretable intermediate results such as gradients, corners, contours, frequency components, and motion vectors.

\subsection{CPU-Based Processing}

The VisionForge architecture demonstrates that a substantial computer vision pipeline can be executed using conventional CPU-based processing without requiring a deep-learning accelerator.

%================================================
% 13. SOFTWARE ENGINEERING DESIGN
%================================================

\section{Software Engineering Design}

The system follows a modular software architecture.

Each processing stage can be implemented as an independent component with defined inputs and outputs.

A simplified structure is:

\begin{verbatim}
VisionForge/
|
+-- preprocessing/
|   +-- grayscale
|   +-- gaussian
|   +-- median
|   +-- clahe
|   +-- fft
|
+-- features/
|   +-- canny
|   +-- log
|   +-- dog
|   +-- harris
|   +-- gabor
|   +-- hog
|
+-- segmentation/
|   +-- kmeans
|   +-- mean_shift
|   +-- region_growing
|   +-- watershed
|
+-- shapes/
|   +-- hough_lines
|   +-- hough_circles
|   +-- contours
|   +-- descriptors
|   +-- pca
|   +-- knn
|
+-- motion/
|   +-- background_subtraction
|   +-- dense_flow
|   +-- sparse_flow
|   +-- affine_motion
|
+-- tests/
+-- main.py
\end{verbatim}

This organization supports independent testing and future extension.

%================================================
% 14. TESTING
%================================================

\section{Testing and Reliability}

Testing was performed across the major processing components.

The implementation included more than 20 tests covering different modules and processing operations.

The testing strategy focused on:

\begin{enumerate}
    \item Verifying successful image loading.
    \item Checking image dimensions.
    \item Validating grayscale conversion.
    \item Testing filtering operations.
    \item Testing feature extraction.
    \item Validating segmentation outputs.
    \item Checking shape detection.
    \item Testing motion-processing components.
    \item Ensuring expected output data types.
    \item Checking for runtime exceptions.
\end{enumerate}

The reported test execution produced zero runtime exceptions during the evaluated test suite.

%================================================
% 15. DISCUSSION
%================================================

\section{Discussion}

VisionForge demonstrates that classical computer vision algorithms can be combined to construct a complete visual processing architecture.

The major advantage of the proposed system is its interpretability. Every stage has a clear mathematical purpose and produces measurable intermediate information.

The preprocessing stage improves the input representation. Feature extraction then converts the image into representations describing edges, corners, gradients, and textures. Segmentation organizes pixels into meaningful groups, while geometric analysis extracts higher-level structural information.

The video module further extends this architecture by introducing temporal information. Background subtraction identifies moving objects, while optical flow provides information about motion between frames.

The benchmark results also demonstrate an important engineering consideration: algorithm selection must account for computational cost in addition to theoretical effectiveness.

%================================================
% 16. LIMITATIONS
%================================================

\section{Limitations}

Although VisionForge provides a broad classical computer vision pipeline, several limitations remain.

\begin{enumerate}
    \item Classical algorithms may be sensitive to illumination and environmental changes.
    \item Parameter selection can strongly influence segmentation and detection results.
    \item K-Means requires the number of clusters to be specified.
    \item Optical flow methods may be affected by large motion, occlusion, or illumination changes.
    \item Classical feature descriptors may be less robust than learned representations for highly complex recognition tasks.
    \item The measured benchmark represents a particular hardware and software environment and therefore should not be interpreted as a universal runtime.
\end{enumerate}

%================================================
% 17. FUTURE WORK
%================================================

\section{Future Work}

Several extensions can be considered for future versions of VisionForge.

\begin{enumerate}
    \item GPU acceleration for computationally expensive operations.
    \item Parallel processing of independent feature extraction modules.
    \item Automatic parameter optimization.
    \item Real-time video processing.
    \item Integration of classical features with machine learning models.
    \item Addition of a graphical user interface.
    \item Automated benchmark visualization.
    \item More advanced temporal scene analysis.
\end{enumerate}

%================================================
% 18. CONCLUSION
%================================================

\section{Conclusion}

VisionForge presents a modular classical computer vision architecture covering the complete processing chain from low-level image preprocessing to video scene dynamics.

The system integrates grayscale conversion, filtering, contrast enhancement, frequency-domain analysis, edge and corner detection, scale-space processing, texture analysis, HOG descriptors, segmentation, geometric shape detection, contour analysis, PCA, KNN, background subtraction, and optical flow.

For a $512\times512$ RGB input image, the measured static processing pipeline required approximately 8.3015 seconds. The system generated 364 Harris corners, 142,884 HOG descriptors, 23 Hough lines, 31 circles, and 28 closed contours.

The empirical results demonstrate that classical computer vision remains useful for understanding the fundamental mechanisms of visual processing. The modular architecture also provides a practical engineering framework in which algorithms can be independently evaluated, profiled, tested, and replaced.

Overall, VisionForge demonstrates how mathematically grounded classical computer vision techniques can be integrated into a complete and measurable visual processing system.

%================================================
% REFERENCES
%================================================

\begin{thebibliography}{99}

\bibitem{canny1986}
J. Canny,
``A Computational Approach to Edge Detection,''
\textit{IEEE Transactions on Pattern Analysis and Machine Intelligence},
vol. PAMI-8, no. 6, pp. 679--698, 1986.

\bibitem{harris1988}
C. Harris and M. Stephens,
``A Combined Corner and Edge Detector,''
in \textit{Proceedings of the Alvey Vision Conference},
pp. 147--151, 1988.

\bibitem{otsu1979}
N. Otsu,
``A Threshold Selection Method from Gray-Level Histograms,''
\textit{IEEE Transactions on Systems, Man, and Cybernetics},
vol. 9, no. 1, pp. 62--66, 1979.

\bibitem{dalal2005}
N. Dalal and B. Triggs,
``Histograms of Oriented Gradients for Human Detection,''
in \textit{Proceedings of the IEEE Conference on Computer Vision and Pattern Recognition},
vol. 1, pp. 886--893, 2005.

\bibitem{lucas1981}
B. D. Lucas and T. Kanade,
``An Iterative Image Registration Technique with an Application to Stereo Vision,''
in \textit{Proceedings of the International Joint Conference on Artificial Intelligence},
pp. 674--679, 1981.

\bibitem{comaniciu2002}
D. Comaniciu and P. Meer,
``Mean Shift: A Robust Approach Toward Feature Space Analysis,''
\textit{IEEE Transactions on Pattern Analysis and Machine Intelligence},
vol. 24, no. 5, pp. 603--619, 2002.

\bibitem{vincent1991}
L. Vincent and P. Soille,
``Watersheds in Digital Spaces: An Efficient Algorithm Based on Immersion Simulations,''
\textit{IEEE Transactions on Pattern Analysis and Machine Intelligence},
vol. 13, no. 6, pp. 583--598, 1991.

\bibitem{duda1972}
R. O. Duda and P. E. Hart,
``Use of the Hough Transformation to Detect Lines and Curves in Pictures,''
\textit{Communications of the ACM},
vol. 15, no. 1, pp. 11--15, 1972.

\end{thebibliography}

%================================================
% END DOCUMENT
%================================================

\end{document}